\documentclass[11pt,a4paper]{report}
\usepackage[utf8]{inputenc}
\usepackage[T1]{fontenc}
\usepackage{amsmath}
\usepackage{graphicx}
\usepackage{hyperref}
\usepackage{geometry}

\geometry{margin=1in}

\title{\textbf{HEP Agent Simulation Framework}\\ \large Model Documentation}
%\author{Development Team}
\date{\today}

\begin{document}

\maketitle
\tableofcontents

\chapter{Introduction}
This document describes the mathematical and logical model underlying the Human Ecology Project (HEP) Agent Simulation.

\chapter{Environment Model}
\section{Grid Structure}
Description of the spatial grid, resolution, and coordinate system.

\section{Habitat Suitability (HEP)}
Description of the HEP index and how it influences agent survival and movement.

\chapter{Agent Dynamics}
\section{Population Structure}
Agents are defined by a set of intrinsic and extrinsic attributes:
\begin{itemize}
    \item \textbf{Age}: Tracked in ticks (1 tick $\approx$ 1 week).
    \item \textbf{Gender}: Binary classification (Male/Female).
    \item \textbf{Resources}: Energy levels accumulated from grid cells.
\end{itemize}

\section{Movement Rules}
Logic governing agent migration and dispersal.

\section{Demographics}
The demographics model manages population growth and mortality through a combination of individual agent rules and macroscopic feedback loops.

\subsection{Birth Model (Verhulst Pressure)}
To ensure population stability and prevent exponential runaway, the simulation utilizes a \textbf{macroscopic fertility scaling factor} (\texttt{K\_fertility}). 

The target birth rate $\phi_{target}$ is calculated using the Verhulst logistic equation:
\begin{equation}
    \phi_{target} = r \cdot \left( 1 - \frac{N}{N_C} \right)
\end{equation}
Where:
\begin{itemize}
    \item $r$: Intrinsic growth rate.
    \item $N$: Current total population count.
    \item $N_C$: Carrying capacity of the environment.
\end{itemize}

The actual fertility of each female agent is then scaled by $K_{fertility}$ to match this target, bounded by \texttt{Kmin} and \texttt{Kmax} to prevent demographic collapse or excessive volatility.

\subsection{Death Model}
Mortality is driven by multiple factors:
\begin{itemize}
    \item \textbf{Natural Mortality}: Age-dependent risk.
    \item \textbf{Starvation}: Triggered when an agent's resources fall below a critical threshold.
    \item \textbf{Environmental Risk}: Probability of death based on HEP suitability (e.g., drowning in water cells).
\end{itemize}

\chapter{Interaction Model}
\section{Resource Competition}
How agents compete for limited resources in a grid cell.

\end{document}
